\begin{textsample}{POS Dim 1 – 2023 – Score 19.00 – 2023/t003440.txt}  \label{ex:f1_pos_064}
\textbf{The} most important question these days is how can we speed up \textbf{the} solutions to \textbf{the} climate crisis?
I ' m convinced we are going to solve \textbf{the} climate crisis.
We ' ve got this.
But \textbf{the} question remains, will we solve it in time?
Others have said we ' re kind of in a race.
I ' ll give you \textbf{the} \textbf{the} shortest definition of \textbf{the} problem.
If I was going to give a one-slide \textbf{slide} show, it would be this \textbf{slide}.
That ' s \textbf{the} troposphere, \textbf{the} lowest part of \textbf{the} atmosphere.
And you already know why it ' s \textbf{blue}.
That ' s \textbf{the} oxygen that refracts \textbf{the} \textbf{blue} light.
And if you could drive a car at \textbf{highway} speed straight up in \textbf{the} air, you ' d get to \textbf{the} top of that \textbf{blue} line in about five to seven minutes.
You could walk it in an hour.
And all of \textbf{the} \textbf{greenhouse} gas pollution is below you.
That ' s what we ' re using as an open sewer.
That ' s \textbf{the} problem.
And it ' s causing a lot of second and third order consequences.
And we saw some of them in \textbf{the} northern-tier cities, including Detroit.
This one ' s from New York City, all \textbf{the} fires in Canada.
And we have gotten used to \textbf{the} fact that \textbf{the} world suffers deep \textbf{droughts} and huge rain bombs and downpours and floods simultaneously. \textbf{The} really ingenious new gravity-measuring satellite has given us, for \textbf{the} first time, \textbf{the} opportunity to see how this plays out worldwide.
We get these huge surpluses of water, \textbf{the} rain bombs and \textbf{the} \textbf{drought} simultaneously.
And as you can see, \textbf{the} amplitude is increasing.
And at both ends of our \textbf{planet} right now, we ' re seeing signs of distress.
Of course, you know, I ' ve often said, every night on \textbf{the} TV news is like a nature hike through \textbf{the} Book of Revelation.
And just today, big flooding in Montpelier, Vermont, in southern Japan, in India.
And I haven ' t done a complete scan, but every single day, it ' s like that.
But in Antarctica, some scientists who are normally pretty levelheaded are getting a little bit freaked out, I would say, is a fair definition, about \textbf{the} lowest level ever, at this point in \textbf{the} year, of \textbf{sea} ice.
And at \textbf{the} other pole, in \textbf{the} North Atlantic, we ' re seeing literally off-the-charts temperatures.
So obviously, \textbf{the} crisis has to be addressed.
And \textbf{the} good news is, as others have often said, we are seeing tremendous progress.
And it starts with \textbf{the} Inflation Reduction Act here in \textbf{the} United States.
President Biden, \textbf{the} Congress, they have passed \textbf{the} best, biggest climate legislation in all of history, and it ' s said to be 369 billion.
But \textbf{the} heavy lifting is done by tax credits, and most of them are open-ended.
And \textbf{the} early applications already show it ' s going to be well over a trillion dollars, maybe 1. 2 trillion.
So this is really good news.
And one month after that passed, Australia changed its government and started changing its laws.
And it ' s now a pro-climate nation.
A month after that, Brazil did \textbf{the} same thing, new president, new policies protecting \textbf{the} Amazon.
And throughout, \textbf{the} European Union has resisted \textbf{the} efforts of Russia to blackmail it into supporting its sadistic and cruel invasion of Ukraine.
And there are other signs of success as well.
China ' s reached its renewables target five years ahead of time.
It ' s still building a lot of \textbf{coal}, but only 50 percent capacity utilization.
So there is lots of good news.
But still, in spite of this progress, \textbf{the} \textbf{emissions} are still going up.
And \textbf{the} crisis is still getting worse faster than we are deploying \textbf{the} solutions.
So maybe it ' s time to look at \textbf{the} obstacles that are standing in our way.
I ' m going to focus on two of them this afternoon.
First of all, \textbf{the} unrelenting opposition from \textbf{the} \textbf{fossil} fuel industry.
A lot of people think they ' re on side and trying to help.
But let me tell you, and \textbf{the} activists will all tell you this, every piece of legislation, whether it ' s at \textbf{the} municipal level, \textbf{the} regional or provincial level, \textbf{the} national level or \textbf{the} international level, they ' re in there with their lobbyists and with their fixers and with their revolving-door colleagues doing everything they can to slow down progress.
So speeding up progress means doing something about this.
They have used fraud on a massive scale.
They ' ve used falsehoods on an industrial scale.
And they ' ve used their legacy political and economic networks, lavishly funded, to capture \textbf{the} policy-making process in too many countries around \textbf{the} world.
And \textbf{the} [ UN ] Secretary General said \textbf{the} \textbf{fossil} fuel industry is \textbf{the} polluted heart of \textbf{the} climate crisis.
Now, that ' s not to say that \textbf{the} men and women who ' ve worked in \textbf{fossil} fuel for \textbf{the} last century and a half are not due our gratitude.
They are, and they didn ' t cause this.
But for decades now, \textbf{the} companies have had \textbf{the} evidence.
They know \textbf{the} truth, and they consciously decided to lie to publics all around \textbf{the} world in order to calm down \textbf{the} political momentum for doing something about it so they can make more money.
It ' s simple as that.
And now they have brazenly seized control of \textbf{the} COP process, especially this year ' s COP in Dubai in \textbf{the} United Arab Emirates.
And concern has been building about this for quite some time.
I remember when there were so many \textbf{fossil} fuel delegates in Madrid, but by \textbf{the} time we got to Glasgow a year and a half ago, \textbf{the} delegates from \textbf{the} \textbf{fossil} fuel companies made up a larger group than \textbf{the} largest national delegation.
And why?
Why?
Because they ' re helping?
They ' re not helping, they ' re trying to stop progress.
And last year in Egypt, they had more delegates than \textbf{the} combined delegations of \textbf{the} ten most affected countries by climate.
And now this year ' s host, which is a petro state, has appointed \textbf{the} president of COP28, in spite of \textbf{the} fact that he has a blatant conflict of interest, he ' s \textbf{the} CEO of \textbf{the} Abu Dhabi National Oil Company owned by Abu Dhabi.
Their \textbf{emissions} are larger than those of ExxonMobil, and they have no credible plan whatsoever to reduce them.
So this is \textbf{the} person in charge of \textbf{the} COP.
He ' s a nice guy.
He ' s a smart guy.
But a conflict of interest is a conflict of interest.
And a matter of fact, they have a plan now to have a new increase in their \textbf{emissions}.
Their plan is to increase \textbf{the} production of both \textbf{oil} and gas by as much as 50 percent by 2030.
Which is \textbf{the} same time frame when \textbf{the} world is trying to reduce \textbf{emissions}, by 50 percent by 2030.
And \textbf{the} same person has been put in charge of both of those efforts.
Direct conflict of interest.
I think it ' s time to say, wait a minute, do you take us for fools?
Do you think you can just completely remove \textbf{the} disguise and we won ' t notice? \textbf{The} \textbf{fossil} fuel industry has captured this process and is slowing it down.
And we need to do something about it.
Now, they plan to increase their \textbf{emissions} -- Getting all hot and bothered here. \textbf{The} head of \textbf{the} International Energy Agency was talking about several \textbf{oil} companies who have publicly pledged to increase their \textbf{oil} and gas and simultaneously pledged to comply with \textbf{the} Paris Agreement.
Do you take us for fools?
You cannot do both of those things simultaneously.
And if you think we believe you when you say that, we beg to differ.
Last year at \textbf{the} COP, \textbf{the} \textbf{fossil} fuel petro states vetoed any mention of phasing down \textbf{fossil} fuels.
They say, oh, that ' s not \textbf{the} problem.
One of \textbf{the} secretary of state equivalent in Saudi Arabia said, we don ' t see this as a discussion about \textbf{fossil} fuels.
Let me tell you, \textbf{the} climate crisis is a \textbf{fossil} fuel crisis. \textbf{The} solutions are going to come from a discussion and collaboration about phasing out \textbf{fossil} fuels.
And there ' s only so much longer they can hold this up and tie us down and keep us from doing \textbf{the} right thing.
Now \textbf{the} director general of this year ' s COP, working under Sultan Al Jaber, he worked for ADNOC also, he was seconded for several years to ExxonMobil.
He said, no, that ' s all changed.
This year, they ' ve changed. \textbf{The} \textbf{fossil} fuel companies are really engaging us.
They have changed.
Yes, they have changed.
They ' ve changed for \textbf{the} worse this year.
You don ' t believe me?
Look at this, BP.
They said,"We ' re going to net zero, there ' s no turning back."A few months ago this year, they turned back and they decided to roll back their investments.
ExxonMobil.
Oh, my God.
This is minor compared to their industrial-scale lying, you know.
But they ' ve had all \textbf{the} TV ads on, particularly in \textbf{the} US, about revolutionizing biofuels."You wouldn ' t believe what algae can do."Well, for 14 years they ran this program.
And they spent 50 percent as much on \textbf{the} ads about \textbf{the} program as they did on actually trying to come up with new biofuels.
And then a couple of months ago, they said,"Oh, we ' ve changed our mind.
We ' re just going to cancel that program."\textbf{The} industry as a whole has not been acting in good faith.
Shell reversed its commitment to increase their investments in \textbf{fossil} fuels.
They announced just a couple of months ago they ' re going to plow that money into expanding \textbf{oil} and gas production instead.
So did they miss \textbf{the} memo?
Well, no, they say they ' ve discovered a way out.
They ' re telling us now that \textbf{the} problem is not \textbf{fossil} fuels.
It ' s \textbf{the} \textbf{emissions} from \textbf{fossil} fuels.
And so all we have to do is to capture \textbf{the} \textbf{emissions} from \textbf{the} \textbf{fossil} fuels.
Now, I ' m all for research and development into trying to capture \textbf{emissions} or sucking it out of \textbf{the} air if they want to do that.
But let ' s don ' t pretend it ' s for real.
Maybe someday it will be.
Maybe someday it will be.
This is \textbf{the} head of ExxonMobil saying,"Not \textbf{oil} and gas, it ' s \textbf{the} \textbf{emissions}."And Abu Dhabi, they have said they ' ve already decreased their \textbf{emissions} from producing \textbf{oil} and gas, 99. 2 percent. \textbf{The} point two makes it sound even more authoritative.
Is that feasible?
Have they really done that?
Is their \textbf{carbon} capture technology that ' s that good?
Well, Climate TRACE can measure.
Now, let ' s just look.
There ' s their CCS capability.
Oh, but they ' ve got a great improvement in \textbf{the} work.
Seven years from now, they ' re going to have a 500 percent increase in their capacity.
But here are \textbf{the} \textbf{emissions} from \textbf{oil} and gas production alone.
And here are \textbf{the} \textbf{emissions} of \textbf{greenhouse} gases from Abu Dhabi, from \textbf{the} UAE.
Now, \textbf{the} pathetic little sliver there compared to \textbf{the} reality, again, do they think we don ' t see what they ' re doing, don ' t understand what they ' re doing?
A lot of people maybe just want to look away from it.
Pretty ironic here in this country.
You know, \textbf{the} coal-burning utilities have been given a new mandate this year by \textbf{the} EPA.
You ' ve got to clean up your act.
Now, you can keep on burning \textbf{coal} and you can keep on producing electricity with it if you use \textbf{carbon} capture and sequestration.
Well, they scream bloody murder.
Another story about it this morning, but they say, that technology?
That ' s not feasible.
It ' s not technically feasible, not economically feasible.
So they need to get their story straight, in my opinion.
So \textbf{the} technology of CCS has been around for a long time, 50 years, and we get used to technologies automatically getting cheaper and better, you know, with \textbf{the} \textbf{computer} chips and \textbf{the} cell phones and all that.
Well, University of Oxford did a whole study, and some come down in cost fast, some slow.
A few are in a category called non-improving technologies.
And that ' s where \textbf{carbon} capture and storage is. 50 years, not any price decline.
Well, could we get a breakthrough in spite of that?
Maybe.
We ' re throwing a lot of money at it.
We had to do that to get \textbf{the} compromise that passed \textbf{the} IRA, and we ought to continue \textbf{the} research.
But again, let ' s don ' t pretend that it ' s here.
It ' s not here now.
And to use it as an excuse, and now they ' ve got another one.
Now it ' s called direct air capture.
These are giant vacuum cleaners that use an awful lot of energy.
It ' s technically feasible, but it ' s extremely expensive.
And it also uses so much energy.
But it is tremendously useful.
You know what it ' s used for? \textbf{The} CEO of one of \textbf{the} largest \textbf{oil} companies in \textbf{the} US had told us what it ' s useful for.
It ' s useful to give them an excuse for not ever stopping \textbf{oil}.
A few months later, she said,"This gives us, this new technology, a license to continue operating."Some people call that a moral hazard.
For them, \textbf{the} moral hazard is not a bug, it ' s a feature, as \textbf{the} old saying goes.
And so let ' s just look, this is state-of-the-art.
It looks pretty impressive, doesn ' t it?
This is \textbf{the} backside of one of these machines.
I had \textbf{the} same thought.
Now this, they ' re improving this.
They ' re improving this.
And \textbf{the} new model, seven years from now, each of these machines is going to be able to capture 27 seconds ' worth of annual \textbf{emissions}.
Woohoo.
That gives them a license to continue producing more and more \textbf{oil} and gas.
Or so they claim.
Well, they also say, \textbf{the} experts, that in order to be economically feasible, it has to come in at 100 dollars a ton or less. \textbf{The} leading company says that 27 years from now, it may come down below three times that amount.
So again, if this sounds kind of incredible, it is not credible.
But they ' re using it in order to gaslight us, literally.
But I didn ' t mean it as a pun, but that ' s what they ' re using it to do.
And we can ' t fall for it. \textbf{The} biggest obstacle to direct air capture is probably physics.
And what I mean by that, they ' ve explained this to me.
CO2 makes up 0. 035 percent of \textbf{the} air so we ' re going to vacuum \textbf{the} other 99. 96 percent to get that little bit out.
Come on, really?
Come on.
And it doesn ' t even pretend to catch \textbf{the} methane or \textbf{the} soot or \textbf{the} mercury or \textbf{the} particulate pollution that kills 9 million people a year around \textbf{the} world.
And that ' s why \textbf{the} climate justice advocates are so down on this kind of nonsense.
They say, here ' s one of them says, look, we need to fight climate change, but don ' t do it by pretending to do it and continue dumping all this pollution on \textbf{the} front line communities that are downwind from \textbf{the} smokestacks.
And that ' s why they have not ever really liked \textbf{the} \textbf{carbon} capture.
And what about all \textbf{the} energy that ' s needed?
Is it going to come from more \textbf{fossil} fuels?
Well, in case of Occidental, yes.
But could it come from more solar and wind?
Well, if so, why not use \textbf{the} solar and wind to replace \textbf{the} \textbf{fossil} fuel-burning plants that are putting all that stuff up there in \textbf{the} first place?
Doesn ' t that make sense?
Am I missing something?
Anyway, this is \textbf{the} environmental justice advocate who is making \textbf{the} point that you can ' t do it without protecting those who are victimized downwind from \textbf{the} smokestacks, downstream from \textbf{the} pollution.
Now, some of \textbf{the} companies, of course, you say we ' re going to offset our \textbf{emissions}.
And again, offsets can play a role, a small role, five to 10 percent according to \textbf{the} Science Based Targets initiative, \textbf{the} most authoritative source on all of this.
But, you know, Chevron just had its offsets analyzed.
It turned out 93 percent of them were worthless and junk.
And that ' s too often \textbf{the} case.
It ' s not a get-out-of-jail-free card.
And their offsets were only aimed at 10 percent of their \textbf{emissions} in \textbf{the} first place.
Eighty percent of \textbf{the} \textbf{fossil} fuel companies just completely ignore \textbf{the} Scope Three \textbf{emissions}, which is \textbf{the} main part of their pollution problem.
So they say that they have increased by 400 percent \textbf{the} share of their spending on energy that goes to green technologies and \textbf{carbon} capture.
And yes, they have, they ' ve increased it all \textbf{the} way up to four percent.
And it ' s not even as good as that, because if you look at \textbf{the} windfall profits they ' ve been getting, what have they done with that?
Well, they gave it back to \textbf{the} shareholders through dividends and stock buybacks.
So \textbf{the} amount of money \textbf{the} \textbf{fossil} fuel industry is investing in renewables and \textbf{carbon} capture is one percent.
Does that mean they are sincere actors, working in good faith?
I don ' t think so.
A lot of people still think so.
They think they ' re on our side.
I don ' t think it ' s in their nature to be on our side.
I think they ' re driven by incentives that push them in \textbf{the} opposite direction.
But in any case, they produce no solutions whatsoever that are scalable or remotely feasible.
And again, they have actively fought against \textbf{the} solutions that others have been trying to bring.
I ' ve seen it personally.
Many of you have seen it personally.
And it ' s happening all over \textbf{the} world.
No progress.
They ' ve even gone \textbf{backwards}.
Now, if a company says, look, we shouldn ' t be excluded from \textbf{the} COP process, we know a lot about energy.
Here ' s a simple test.
Do they actually have, does \textbf{the} company have real net-zero commitments, a genuine phase-down plan, yes or no?
Are they committed to full disclosure?
Yes or no?
Are they going to spend windfall profits on \textbf{the} transition?
Yes or no?
Are they committed to transparency?
Will they end their anti-climate lobbying?
Will they end their greenwashing?
Will they be in favor of reforming \textbf{the} COP process so that \textbf{the} petro states don ' t have an absolute veto on anything \textbf{the} world wants to discuss or act on?
We ' ve got to change that.
In order to move faster, we have got to empower \textbf{the} global community in a way that frees them from \textbf{the} hammerlock that \textbf{the} \textbf{fossil} fuel companies have on them today.
And any company that doesn ' t pass this test, I ' m telling you, ought to be prohibited from taking part in \textbf{the} COP process.
What about \textbf{the} company that ' s led by \textbf{the} president of this year ' s COP, \textbf{the} Abu Dhabi National Oil Company?
Well, it ' s rated one of \textbf{the} least responsible of all \textbf{the} \textbf{oil} and gas companies.
It no longer even releases information on any of its \textbf{emissions}, has no transition plan, no short or mid-term target scores, a pathetic three on 100-scale on their transition plans.
And \textbf{the} nation ' s plan is rated highly insufficient.
They just put out a new one, it ' s still insufficient.
So that ' s \textbf{the} first obstacle.
Second one ' s a lot clearer.
This one has to be removed in order to get to \textbf{the} faster solutions.
But in order to remove it, we ' ve got to do something about \textbf{the} second obstacle, and that is \textbf{the} financial system, \textbf{the} global allocation of capital and \textbf{the} subsidies for \textbf{fossil} fuels.
Governments last year around \textbf{the} world subsidized with taxpayer money more than a trillion dollars, subsidizing \textbf{fossil} fuels.
That ' s five times larger than \textbf{the} amount in 2020.
Are we going in \textbf{the} right direction or \textbf{the} wrong direction?
Well, there are a lot of good things I ' ll show you, but this has to stop.
And \textbf{the} 60 largest global banks have put 5. 5 trillion extra dollars into \textbf{the} \textbf{fossil} fuel companies since \textbf{the} Paris agreement.
And preposterously, 49 of them have also signed net-zero pledges.
So we...
Global allocation of capital. \textbf{The} borrowing of private capital is ruled out for many developing countries.
Nigeria has to pay an interest rate seven times higher than Europe or \textbf{the} US.
Political risk, corruption risk, currency risk, off-take risk, it ' s rule of law risk.
That ' s why \textbf{the} World Bank and \textbf{the} other multilateral development banks are supposed to take those top layers of risk off.
And that ' s why we need reform in \textbf{the} World Bank.
Thank goodness so many fought to get a new head of \textbf{the} World Bank, and I ' m very excited about his progress.
Now, when these obstacles are removed, we can really accelerate progress.
They ' ll still be problems, of course, but we have everything we need and proven deployment models to reduce \textbf{emissions} 50 percent in \textbf{the} next seven years.
But we will still need better grids, more resilient grids, more solar and wind, much more regenerative \textbf{agriculture}.
That ' s a way to really pull \textbf{carbon} out of \textbf{the} air.
More electric vehicles and charging stations and more energy storage.
I ' ll show you that just before we quit here.
And more green hydrogen, so more electrolyzers to produce it.
And these are all surmountable obstacles.
Already, renewables account for 90 percent of all \textbf{the} new electricity generation being installed each year worldwide.
Ninety-three percent solar and wind in India last year.
And of course, \textbf{the} solar ' s taking off and so is wind and so are \textbf{the} electric vehicles and battery storage.
It has grown so dramatically.
You ' ve got to rescale \textbf{the} graph to look at \textbf{the} trillion-dollar industry emerging.
Now, just seven years ago, there was one gigafactory.
I learned from one of you at a session a couple of hours ago and had my staff research it, there are now 195 gigafactories and another 300 in \textbf{the} pipeline.
So will we succeed?
Some people fight with a vulnerability to despair.
You know, \textbf{the} old cliche,"denial ain ' t just a river in Egypt."Despair ain ' t just a tire in \textbf{the} trunk.
There are people that are vulnerable.
But let me close with what I regard as amazingly good news.
What if we could stop \textbf{the} increase in temperatures?
Well, if you look at \textbf{the} temperature increases, if we get to true net-zero, astonishingly, global temperatures will stop going up with a lag time of as little as three to five years.
They used to think that positive feedback loops would keep that process going.
No, it will not. \textbf{The} temperatures will stop going up. \textbf{The} ice will continue melting and some other things will continue, but we can stop \textbf{the} increase of temperatures.
Even better, if we stay at true net zero, in as little as 30 years, half of all \textbf{the} human-caused CO2 will come out of \textbf{the} atmosphere into \textbf{the} upper ocean and \textbf{the} trees and vegetation.
So young people are demanding that we do \textbf{the} right thing.
Do not be vulnerable to despair.
We are going to do this.
And if you doubt that we, as human beings have \textbf{the} will to act, please always remember that \textbf{the} will to act is itself a \textbf{renewable} resource.
Thank you very much.
Thank you.

% matched lemmas: agriculture, backwards, blue, carbon, coal, computer, drought, emission, fossil, greenhouse, highway, oil, planet, renewable, sea, slide, the
\end{textsample}
